\subsubsection{Perron-boundary Adams--M\"obius renormalisation}
\label{subsubsec:perron-boundary-adams-mobius-renormalisation}

This subsection contains the Perron-boundary theorem used in the rest of the
paper.  It does not assume the strict twisted gap: it keeps all peripheral
phases \(\lambda\zeta\), subtracts only the logarithmic part coming from
\(\zeta=1\), and evaluates the constants by reduced determinant finite parts.
The proof is split into four labelled components, but these components are
not separate theorem spines; they are the primitive-channel, class-Mertens,
determinant, and product steps of the single boundary theorem.

\begin{definition}[Peripheral Peter--Weyl data]
  \label{def:peripheral-peter-weyl-data}
  Assume \(A\) is primitive with Perron root \(\lambda>1\).  For
  \(\pi\in\Irr(G)\), define
  \[
    \Gamma_\pi:=
    \{\zeta\in\TT:\lambda\zeta\in\spec(A_\pi)\}.
  \]
  For \(\zeta\in\Gamma_\pi\), let
  \(m_{\pi,\zeta}\) be the algebraic multiplicity of
  \(\lambda\zeta\) in \(A_\pi\); equivalently,
  \(m_{\pi,\zeta}\) is the rank, or trace, of the Riesz spectral
  projection onto the generalized \(\lambda\zeta\)-eigenspace.  The
  optional semisimplicity certificate below further identifies this
  generalized eigenspace with the ordinary eigenspace.  Put
  \[
    s_\pi:=m_{\pi,1},
  \]
  with \(s_\pi=0\) if \(1\notin\Gamma_\pi\).  Let
  \[
    \Lambda_\perp:=
    \max\{|\alpha|:\alpha\in\spec(A_\pi),\ \pi\in\Irr(G),\
      |\alpha|<\lambda\},
  \]
  after deleting all eigenvalues \(\lambda\zeta\) with
  \(\zeta\in\Gamma_\pi\).  If the remaining set is empty, set
  \(\Lambda_\perp=0\).  Finally set
  \[
    \Lambda_\partial:=\max\{\Lambda_\perp,\lambda^{1/2}\}.
  \]
  The term \(\lambda^{1/2}\) is the same Adams--M\"obius divisor-loss scale
  that appears when the \(m>1\) divisors are bounded in
  \Cref{thm:primitive-peter-weyl-trace}.

  The spectral entries \(\Gamma_\pi\) and \(m_{\pi,\zeta}\) are not by
  themselves a branch convention for boundary logarithms.  Write
  \(P_\pi(z)=\det(I-zA_\pi)\), let \(\Log_0 P_\pi(z)^{-1}\) denote the
  determinant-logarithm germ normalized to be \(0\) at \(z=0\), and put
  \(s_\pi=m_{\pi,1}\), as above.  Let \(e_G\) be the exponent of \(G\).
  The open-disc primitive coefficients are determined by the smaller finite
  datum
  \[
    (G,\operatorname{Pow}_{G},\{P_\pi(z)=\det(I-zA_\pi)\}_\pi),
    \qquad
    \operatorname{Pow}_{G}:=\{C\mapsto C^a:0\le a<e_G\}.
  \]
  The identity \(g^{m+e_G}=g^m\) makes this finite table sufficient for all
  Adams coefficients \(a_{\rho,\pi}^{(m)}\).
  This means the coefficients of the primitive power series inside
  \(|z|<\lambda^{-1}\); it is not a branch-normalized finite-part datum at
  \(z=\lambda^{-1}\).
  The Perron-boundary constants require the following larger but still finite
  boundary datum.  For each \(\pi\), factor
  \begin{equation}\label{eq:perron-boundary-factorization-local}
    P_\pi(t/\lambda)
    =
    \prod_{\zeta\in\Gamma_\pi}(1-t\zeta)^{m_{\pi,\zeta}}
      Q_{\partial,\pi}(t),
    \qquad Q_{\partial,\pi}(1)\ne0 .
  \end{equation}
  For \(0\le t<1\), let
  \(\Log_{\pi,\zeta}(1-t\zeta)\) be the unique continuous branch along this
  segment with value \(0\) at \(t=0\).  Since \(Q_{\partial,\pi}(0)=1\),
  let \(\Log_{\partial,\pi}Q_{\partial,\pi}(1)\) denote the value obtained by
  continuing the logarithm of \(Q_{\partial,\pi}(t)\), normalized to be \(0\)
  at \(t=0\), along \(0\le t<1\).  These branch functions, not a list of
  determinant finite parts, are the boundary branch input.  Product constants
  use no additional branch datum: after
  \Cref{lem:renormalized-reality-positivity} their positive values are read
  with the ordinary real logarithm on \(\mathbb R_{>0}\).
  \begin{equation}\label{eq:perron-boundary-renormalized-data-local}
  \begin{aligned}
    D_{\partial}^{\mathrm{ren}}(A,\tau;G)
    :=
    \bigl(&G,\operatorname{Pow}_{G},\{P_\pi\}_{\pi\in\Irr(G)},\lambda,\\
      &\{\Gamma_\pi\}_{\pi\in\Irr(G)},
        \{m_{\pi,\zeta}\}_{\zeta\in\Gamma_\pi,\ \pi\in\Irr(G)},
        \{\Log_{\pi,\zeta}\}_{\zeta\in\Gamma_\pi,\ \pi\in\Irr(G)},
        \{\Log_{\partial,\pi}\}_{\pi\in\Irr(G)}\bigr).
  \end{aligned}
  \end{equation}
  Here \(\operatorname{Pow}_{G}\) denotes the finite conjugacy-class power
  table modulo \(e_G\) used to compute the Adams coefficients
  \(a_{\rho,\pi}^{(m)}\).  The polynomial \(Q_{\partial,\pi}\) is determined
  by the displayed factorisation, so it need not be a separate coordinate of
  \(D_{\partial}^{\mathrm{ren}}\); it is named to define the second family of
  logarithm branches.  The multiplicity
  \(s_\pi=m_{\pi,1}\), the spectral radii
  \(\Lambda_\perp,\Lambda_\partial\), and optional semisimplicity certificates
  are finite hypotheses or verification data for the estimates below, but they
  are not boundary logarithm inputs.

  The tuple \(D_{\partial}^{\mathrm{ren}}(A,\tau;G)\) is finite input data;
  no infinite sequence of determinant finite parts is included.  The
  determinant finite parts, fixed-label coordinates, and product constants are
  outputs computed from it.  Once this finite datum is fixed, define for every
  \(\pi\in\Irr(G)\) and \(k\ge1\) the derived branch-normalized determinant
  layer directly by the Abel path
  \begin{equation}\label{eq:perron-boundary-renormalized-b-data}
    B_{k,\pi}^{\mathrm{ren}}
    =\mathcal B_{k,\pi}^{\mathrm{ren}}:=
    \begin{cases}
      \displaystyle
      \lim_{t\uparrow1}
      \left(
        \Log_0\det(I-t\lambda^{-1}A_\pi)^{-1}
        +s_\pi\log(1-t)
      \right), & k=1,\\[1.2em]
      \displaystyle
      \Log_0\det(I-\lambda^{-k}A_\pi)^{-1}, & k\ge2.
    \end{cases}
  \end{equation}
  Equivalently, from \eqref{eq:perron-boundary-factorization-local}, the
  \(k=1\) entry is
  \begin{equation}\label{eq:perron-boundary-renormalized-b1-factor}
    B_{1,\pi}^{\mathrm{ren}}
    =
    -\sum_{\substack{\zeta\in\Gamma_\pi\\\zeta\ne1}}
      m_{\pi,\zeta}\Log_{\pi,\zeta}(1-\zeta)
    -\Log_{\partial,\pi} Q_{\partial,\pi}(1).
  \end{equation}
  Consequently the reduced fixed-label boundary coordinate is the derived
  quantity
  \begin{equation}\label{eq:perron-boundary-fren-derived-data}
    F_\varrho^{\mathrm{ren}}
    =
    \sum_{r,m\ge1}\frac{\mu(m)}{rm}
      \sum_{\pi\in\Irr(G)}
        a_{\varrho,\pi}^{(m)}B_{rm,\pi}^{\mathrm{ren}}.
  \end{equation}
  For \(k=1\), the logarithm is continued only along the Abel path
  \(0\le t<1\); for \(k\ge2\), the evaluation is an ordinary interior-disc
  value.  Thus later ``finite determinant'' statements refer to the finite
  determinant polynomials together with this finite branch-normalized recipe
  and, when an effective approximation is requested, a finite truncation level
  and rational tail majorant chosen after the theorem.  The values
  \(B_{k,\pi}^{\mathrm{ren}}\) are outputs of that recipe, not independent
  inputs.
\end{definition}

\paragraph{Boundary conventions.}
\begin{definition}[Boundary conventions for determinant and product coordinates]
  \label{def:perron-boundary-logarithm-convention}
  Fix the peripheral data of
  \Cref{def:peripheral-peter-weyl-data} and write
  \(P_\pi(z)=\det(I-zA_\pi)\).  All determinant logarithms below are the
  normalized germs \(\Log_0\) at \(z=0\), continued only along the indicated
  Abel path or evaluated at an interior point.  For \(k\ge1\) and
  \(\pi\in\Irr(G)\), the boundary logarithm is the value
  \(\mathcal B_{k,\pi}^{\mathrm{ren}}\) defined in
  \eqref{eq:perron-boundary-renormalized-b-data}.
  Thus every \(k=1\) Perron-boundary block is read as a reduced Abel finite
  part: exactly the \(s_\pi=m_{\pi,1}\) copies of the non-oscillatory factor
  \(1-t\) are cancelled, while the remaining peripheral factors
  \(1-t\zeta\), \(\zeta\ne1\), are evaluated on the Abel branch reached from
  the origin.  Every \(k\ge2\) term is an ordinary interior-disc logarithm.
  This definition is the only branch convention used for the boundary product
  coordinate package.  Later phrases such as ``branch-compatible'',
  ``reduced determinant layer'', ``strict-gap determinant value'', and
  ``deleted scalar triple'' refer back to the present rule and do not add new
  choices.  The scalar/product logarithm convention is also fixed here, but it
  is not an independent boundary datum: after
  \Cref{lem:renormalized-reality-positivity} the product constants are
  positive real numbers, and their logarithms are the ordinary real logarithm
  on \(\mathbb R_{>0}\).

  In the strict-gap specialization, where
  \(s_{\mathbf1}=1\), \(s_\pi=0\) for \(\pi\ne\mathbf1\), and no non-trivial
  block has a Perron-boundary eigenvalue, the singular scalar boundary value
  \(\mathcal B_{1,\mathbf1}^{\mathrm{ren}}\) is not used in non-trivial
  fixed-label coordinates.  For those coordinates we write
  \(\mathcal L_{k,\pi}:=\mathcal B_{k,\pi}^{\mathrm{ren}}\) except that the
  triple \((k,\pi)=(1,\mathbf1)\), equivalently
  \((r,m,\pi)=(1,1,\mathbf1)\), is deleted before any endpoint value is
  chosen.  Non-trivial \(k=1\) terms are Abel boundary values and all
  \(k\ge2\) terms, including scalar terms created by Adams powers, remain
  ordinary interior determinant logarithms.

  Effective finite enclosures use the same convention with a finite
  truncation level and a rational tail majorant chosen after the analytic
  theorem has been proved.  Those truncation and interval data are
  reproducibility outputs; they are not extra coordinates in
  \(D_{\partial}^{\mathrm{ren}}\) and they do not modify the branch rule.
\end{definition}

\begin{lemma}[Uniform interior determinant layers]
  \label{lem:perron-boundary-uniform-interior-layers}
  Fix the finite boundary datum
  \(D_{\partial}^{\mathrm{ren}}\) of
  \eqref{eq:perron-boundary-renormalized-data-local}.  For every
  \(k\ge2\) and every \(\pi\in\Irr(G)\), the point
  \(\lambda^{-k}\) lies in the common open disc on which the normalized
  determinant logarithm is the absolutely convergent trace series, and
  \[
    B_{k,\pi}^{\mathrm{ren}}
    =
    \Log_0 P_\pi(\lambda^{-k})^{-1}
    =
    \sum_{\ell\ge1}
      \frac{\lambda^{-k\ell}}{\ell}\Tr(A_\pi^\ell).
  \]
  Moreover there is a constant \(C\), depending only on the finite
  skew-product matrix and not on \(k\) or \(\pi\), such that
  \begin{equation}\label{eq:perron-boundary-uniform-interior-layer-bound}
    |B_{k,\pi}^{\mathrm{ren}}|\le C\lambda^{1-k}
    \qquad(k\ge2,\ \pi\in\Irr(G)).
  \end{equation}
  Hence all \(k\ge2\) determinant layers are derived uniformly from the
  finite polynomials \(P_\pi\), the Perron root \(\lambda\), and the
  normalized germ \(\Log_0\); no peripheral factorization, finite-part
  subtraction, or extra infinite list of \(B\)-values is involved.
\end{lemma}

\begin{proof}
  The Peter--Weyl decomposition of the finite skew-product matrix gives
  \(\rho(A_\pi)\le\lambda\) for every \(\pi\).  Since \(\Irr(G)\) is finite,
  Jordan normal form for the finitely many blocks gives a constant \(C_0\)
  such that
  \[
    |\Tr(A_\pi^\ell)|\le C_0\lambda^\ell
    \qquad(\ell\ge1,\ \pi\in\Irr(G)).
  \]
  For \(k\ge2\),
  \(\rho(\lambda^{-k}A_\pi)\le\lambda^{1-k}\le\lambda^{-1}<1\).
  Therefore the logarithm branch normalized at the origin is represented by
  the ordinary determinant trace expansion at \(z=\lambda^{-k}\):
  \[
    \Log_0\det(I-\lambda^{-k}A_\pi)^{-1}
    =
    \sum_{\ell\ge1}\frac{\lambda^{-k\ell}}{\ell}\Tr(A_\pi^\ell).
  \]
  The trace bound gives
  \[
    |B_{k,\pi}^{\mathrm{ren}}|
    \le
    C_0\sum_{\ell\ge1}\frac{\lambda^{(1-k)\ell}}{\ell}
    \le C\lambda^{1-k},
  \]
  uniformly in \(k\ge2\) and \(\pi\), after increasing \(C\) once for
  the fixed number \(\lambda>1\).  This proves both the formula and the
  uniform bound.
\end{proof}

\begin{lemma}[Finite-cover peripheral expansion and optional certificate]
  \label{lem:finite-cover-peripheral-expansion}
  Let \(\widetilde A\) be the finite skew-product adjacency matrix attached
  to \(\tau\).  Suppose the peripheral algebraic data of
  \Cref{def:peripheral-peter-weyl-data} have been computed: for each
  \(\pi\), the list \(\Gamma_\pi\) consists of the roots
  \(\lambda\zeta\) of \(\det(xI-A_\pi)\) on the circle \(|x|=\lambda\),
  \(m_{\pi,\zeta}\) is the algebraic multiplicity of that root, and
  \(\Lambda_\perp<\lambda\) is the maximum modulus of the remaining roots
  (or \(0\) if there are none).  These algebraic data are obtained by the
  finite tests
  \begin{enumerate}[label=\textup{(\roman*)},leftmargin=*]
    \item compute the roots \(\alpha\) of
      \(\det(xI-A_\pi)\) with \(|\alpha|=\lambda\), write
      \(\alpha=\lambda\zeta\), and record
      \(\Gamma_\pi\) and the algebraic multiplicity \(m_{\pi,\zeta}\);
    \item set \(\Lambda_\perp\) to the maximum modulus of all remaining
    roots, or to \(0\) if there are none, and check
    \(\Lambda_\perp<\lambda\).
  \end{enumerate}
  These tests use only determinants and finite algebraic comparison of
  moduli; no limiting orbit counts enter.
  If one also wants the stronger semisimplicity certificate
  \(\mathfrak s_\partial\), one adds the finite rank test
  \[
    \dim\ker(A_\pi-\lambda\zeta I)=m_{\pi,\zeta}
  \]
  for every recorded peripheral root, equivalently that the local factor of
  the minimal polynomial at \(x=\lambda\zeta\) is \(x-\lambda\zeta\).  This
  additional certificate may be checked from the finite skew-product matrix
  \(\widetilde A\) after its Frobenius normal form has no nilpotent coupling
  between same-radius peripheral components.  The determinant and trace
  identities below use the algebraic multiplicities recorded above.

  Under the peripheral algebraic multiplicity data above, for every
  \(\pi\in\Irr(G)\),
  \begin{equation}\label{eq:finite-cover-peripheral-trace-expansion}
    \Tr(A_\pi^n)=
    \lambda^n\sum_{\zeta\in\Gamma_\pi}m_{\pi,\zeta}\zeta^n
    +O(\Lambda_\perp^n),
    \qquad n\ge1,
  \end{equation}
  with a constant depending only on the finite determinant model.  In
  particular all peripheral contributions are finite root-of-unity phases
  inherited from the finite cover, and the trace remainder is strictly below
  modulus \(\lambda\).
\end{lemma}

\begin{proof}
  By \Cref{prop:kernel-peter-weyl-block-diagonalisation},
  \(\widetilde A\) is similar to
  \(\bigoplus_{\pi\in\Irr(G)}A_\pi^{\oplus\dim\pi}\).  Thus the peripheral
  eigenvalues of the twisted blocks are exactly the Peter--Weyl pieces of
  the peripheral spectrum of the finite non-negative matrix
  \(\widetilde A\), counted with multiplicity.  On each irreducible
  component of \(\widetilde A\) with spectral radius \(\lambda\), the
  Perron--Frobenius cyclic theorem gives a finite cyclic peripheral set
  \(\lambda\zeta\), with \(\zeta\) a root of unity.  Components with smaller
  spectral radius contribute only eigenvalues of modulus \(<\lambda\).  The
  optional semisimplicity certificate in the statement would rule out higher
  Jordan blocks for the twisted Peter--Weyl pieces at these peripheral
  values, but the trace identity itself does not require this.

  Fix \(\pi\).  Put \(A_\pi\) in Jordan normal form.  If the optional
  semisimplicity certificate is present, every peripheral Jordan block is
  \(1\times1\).
  Without it, a Jordan block \(J=\lambda\zeta I+N\) of size \(d\) still
  satisfies
  \[
    \Tr(J^n)
    =
    \Tr\left(\sum_{a=0}^{d-1}\binom na(\lambda\zeta)^{n-a}N^a\right)
    =
    d(\lambda\zeta)^n,
  \]
  because \(N^a\) is nilpotent, and hence has trace \(0\), for \(a\ge1\).
  Thus the total contribution of all Jordan blocks at \(\lambda\zeta\) is
  \(m_{\pi,\zeta}(\lambda\zeta)^n\), with \(m_{\pi,\zeta}\) the algebraic
  multiplicity.  Every remaining eigenvalue has modulus at most
  \(\Lambda_\perp\).  Applying the same trace calculation to the finitely
  many non-peripheral Jordan blocks, the sum of the remaining powers is
  \(O(\Lambda_\perp^n)\), with the convention that the term is zero if
  \(\Lambda_\perp=0\).  This proves
  \eqref{eq:finite-cover-peripheral-trace-expansion}.
\end{proof}

\begin{proposition}[Primitive-channel estimate at the Perron boundary]
  \label{prop:peripheral-primitive-channel-asymptotic-component}
  Under the peripheral algebraic multiplicity hypotheses of
  \Cref{def:peripheral-peter-weyl-data,lem:finite-cover-peripheral-expansion},
  every irreducible \(\varrho\in\Irr(G)\) satisfies
  \begin{equation}\label{eq:peripheral-primitive-channel-asymptotic-component}
    c_{n,\varrho}(A,\tau)
    =
    \frac{\lambda^n}{n}
      \sum_{\zeta\in\Gamma_\varrho}m_{\varrho,\zeta}\zeta^n
    +
    O\!\left(\frac{\Lambda_\partial^n}{n}\right).
  \end{equation}
  More explicitly, the \(m=1\) term in the Adams--M\"obius trace formula
  gives the displayed peripheral sum and the interior
  \(O(\Lambda_\perp^n/n)\) remainder, while the full \(m>1\) divisor
  correction is
  \[
    \frac{1}{n}\sum_{\substack{m\mid n\\m>1}}\mu(m)
      \sum_{\pi\in\Irr(G)}
        a_{\varrho,\pi}^{(m)}\Tr(A_\pi^{n/m})
    =
    O\!\left(\frac{\lambda^{n/2}}{n}\right).
  \]
\end{proposition}

\begin{proof}
  Start from the coefficient formula
  \eqref{eq:primitive-peter-weyl-trace}:
  \[
    c_{n,\varrho}(A,\tau)
    =
    \frac{1}{n}\Tr(A_\varrho^n)
    +
    \frac{1}{n}\sum_{\substack{m\mid n\\m>1}}\mu(m)
      \sum_{\pi\in\Irr(G)}
        a_{\varrho,\pi}^{(m)}\Tr(A_\pi^{n/m}).
  \]
  Apply \Cref{lem:finite-cover-peripheral-expansion} to the first trace.
  This gives
  \[
    \frac{1}{n}\Tr(A_\varrho^n)
    =
    \frac{\lambda^n}{n}
      \sum_{\zeta\in\Gamma_\varrho}m_{\varrho,\zeta}\zeta^n
    +O(\Lambda_\perp^n/n).
  \]
  For the remaining divisors, the Adams coefficients are uniformly bounded
  because \(G\) has finitely many irreducible characters, and
  \(\rho(A_\pi)\le\lambda\) by
  \Cref{prop:twisted-spectral-radius-bound}.  Hence
  \(|\Tr(A_\pi^\ell)|\le C\lambda^\ell\), uniformly in \(\pi\) and
  \(\ell\).  If \(m>1\) divides \(n\), then \(n/m\le n/2\), and therefore
  \[
    \sum_{\substack{m\mid n\\m>1}}\lambda^{n/m}
    \le \sum_{k\le n/2}\lambda^k
    =O(\lambda^{n/2}).
  \]
  This proves the divisor-loss estimate and hence
  \eqref{eq:peripheral-primitive-channel-asymptotic-component}.
\end{proof}

\begin{proposition}[Class-Mertens finite part at the Perron boundary]
  \label{prop:renormalized-class-mertens-component}
  Assume the Perron-boundary algebraic multiplicity data above.  For every
  \(\varrho\in\Irr(G)\), define
  \begin{equation}\label{eq:renormalized-primitive-channel-finite-part}
    R_\varrho:=
    \lim_{t\uparrow1}
      \left(\Pi_\varrho(t\lambda^{-1})+s_\varrho\log(1-t)\right).
  \end{equation}
  Then the limit exists and
  \begin{equation}\label{eq:renormalized-primitive-channel-constant}
    R_\varrho=
    -\sum_{\substack{\zeta\in\Gamma_\varrho\\\zeta\ne1}}
      m_{\varrho,\zeta}\Log(1-\zeta)
    +R_{\varrho,\perp},
  \end{equation}
  where \(R_{\varrho,\perp}\) is the absolutely convergent interior and
  Adams-divisor correction obtained after subtracting all peripheral
  terms.  Moreover, for every \(N\ge1\),
  \begin{equation}\label{eq:renormalized-channel-partial-sum}
    \sum_{n\le N}\frac{c_{n,\varrho}(A,\tau)}{\lambda^n}
    =
    s_\varrho H_N
    +R_\varrho
    -
    \sum_{\substack{\zeta\in\Gamma_\varrho\\\zeta\ne1}}
      m_{\varrho,\zeta}\sum_{n>N}\frac{\zeta^n}{n}
    +O(\Lambda_\partial^N/(N\lambda^N)).
  \end{equation}
  In particular the oscillatory tails with \(\zeta\ne1\) are
  \(O(N^{-1})\).

  For a conjugacy class \(C\subset G\), put
  \begin{equation}\label{eq:renormalized-class-density}
    \delta_C:=
    \frac{|C|}{|G|}
      \sum_{\varrho\in\Irr(G)}
        \chi_\varrho^\ast(C)s_\varrho .
  \end{equation}
  Then the pointwise class-count expansion is
  \begin{equation}\label{eq:renormalized-class-pointwise-expansion}
    p_{n,C}
    =
    \frac{|C|}{|G|}
    \frac{\lambda^n}{n}
      \sum_{\varrho\in\Irr(G)}
        \chi_\varrho^\ast(C)
        \sum_{\zeta\in\Gamma_\varrho}
          m_{\varrho,\zeta}\zeta^n
    +O(\Lambda_\partial^n/n),
  \end{equation}
  and the summatory Mertens formula is
  \begin{equation}\label{eq:renormalized-class-mertens-harmonic}
    \sum_{n\le N}\frac{p_{n,C}}{\lambda^n}
    =
    \delta_C H_N+
    \frac{|C|}{|G|}
      \sum_{\varrho\in\Irr(G)}
        \chi_\varrho^\ast(C)R_\varrho
    +O(N^{-1}).
  \end{equation}
  Equivalently,
  \begin{equation}\label{eq:renormalized-class-mertens-log}
    \sum_{n\le N}\frac{p_{n,C}}{\lambda^n}
    =
    \delta_C\log N+M_C^{\mathrm{ren}}+O(N^{-1}),
    \qquad
    M_C^{\mathrm{ren}}:=
    \frac{|C|}{|G|}
      \sum_{\varrho\in\Irr(G)}
        \chi_\varrho^\ast(C)(s_\varrho\gamma+R_\varrho).
  \end{equation}
  The exponent \(\delta_C\) is non-negative, and the displayed
  \(M_C^{\mathrm{ren}}\) is real, by
  \Cref{lem:renormalized-reality-positivity}.
\end{proposition}

\begin{proof}
  This proof uses only the primitive-channel estimate of
  \Cref{prop:peripheral-primitive-channel-asymptotic-component}, not the
  consolidated theorem stated below.  Define
  \begin{equation}\label{eq:renormalized-channel-error-b}
    b_{\varrho,n}:=
    \lambda^{-n}c_{n,\varrho}(A,\tau)
    -\frac1n
      \sum_{\zeta\in\Gamma_\varrho}
        m_{\varrho,\zeta}\zeta^n .
  \end{equation}
  By \Cref{prop:peripheral-primitive-channel-asymptotic-component},
  \[
    |b_{\varrho,n}|
    \le
    C\frac{(\Lambda_\partial/\lambda)^n}{n}.
  \]
  Hence \(\sum_{n\ge1}|b_{\varrho,n}|<\infty\), because
  \(\Lambda_\partial/\lambda<1\).  Therefore, for \(0\le t<1\),
  \[
    \Pi_\varrho(t\lambda^{-1})
    =
    \sum_{\zeta\in\Gamma_\varrho}
      m_{\varrho,\zeta}\sum_{n\ge1}\frac{(\zeta t)^n}{n}
    +\sum_{n\ge1}b_{\varrho,n}t^n,
    \qquad
    \sum_n|b_{\varrho,n}|<\infty .
  \]
  The term \(\zeta=1\) is \(-s_\varrho\log(1-t)\).  For
  \(\zeta\ne1\), the Abel limit
  \(\sum_{n\ge1}\zeta^n t^n/n\to-\Log(1-\zeta)\) exists in the branch
  reached from the origin along \(0\le t<1\).  This proves existence of
  \(R_\varrho\) and the expression
  \eqref{eq:renormalized-primitive-channel-constant}, with the interior
  correction explicitly given by
  \[
    R_{\varrho,\perp}=\sum_{n\ge1}b_{\varrho,n}.
  \]
  The same decomposition of partial sums gives
  \eqref{eq:renormalized-channel-partial-sum}; Dirichlet summation gives
  \(\sum_{n>N}\zeta^n/n=O(N^{-1})\) for \(\zeta\ne1\), and the interior tail
  is geometric with size \(O(\Lambda_\partial^N/(N\lambda^N))\).

  Apply the class decomposition
  \eqref{eq:class-primitive-coefficients} and substitute
  \Cref{prop:peripheral-primitive-channel-asymptotic-component}.  The finite character
  sum gives the pointwise expansion
  \eqref{eq:renormalized-class-pointwise-expansion}; the coefficient of
  the non-oscillatory harmonic term is exactly
  \(\delta_C\).  Summing the channel identities
  \eqref{eq:renormalized-channel-partial-sum} over
  \(\varrho\) with the Hermitian class coefficients gives
  \eqref{eq:renormalized-class-mertens-harmonic}.  Finally
  \(H_N=\log N+\gamma+O(N^{-1})\), which changes the displayed constant to
  \(M_C^{\mathrm{ren}}\) and proves
  \eqref{eq:renormalized-class-mertens-log}.
\end{proof}

\begin{lemma}[Full-boundary majorant for fixed-label determinant limits]
  \label{lem:perron-boundary-fixed-label-majorant}
  Under the peripheral algebraic multiplicity hypotheses of
  \Cref{def:peripheral-peter-weyl-data,lem:finite-cover-peripheral-expansion},
  let \(c_{r,m,\pi}\) be a bounded complex coefficient family indexed by
  \(r,m\ge1\) and \(\pi\in\Irr(G)\).  For \(0<t<1\) put
  \[
    L_{r,m,\pi}(t):=
    \Log_0\det(I-t^{rm}\lambda^{-rm}A_\pi)^{-1}.
  \]
  Then the following statements hold.
  \begin{enumerate}[label=\textup{(\roman*)},leftmargin=*]
    \item The boundary values for the \(rm=1\) terms exist after the
    prescribed finite-part subtraction:
    \[
      \lim_{t\uparrow1}
      \bigl(L_{1,1,\pi}(t)+s_\pi\log(1-t)\bigr)
      =B_{1,\pi}^{\mathrm{ren}}.
    \]
    The value is the factorized expression
    \eqref{eq:perron-boundary-renormalized-b1-factor}.
    \item For \(rm\ge2\), the logarithms are ordinary interior-disc values
    at \(t=1\), and uniformly for \(0\le t\le1\),
    \begin{equation}\label{eq:perron-boundary-fixed-label-majorant-tail}
      |L_{r,m,\pi}(t)|
      \le C\,\lambda^{1-rm}.
    \end{equation}
    Consequently
    \[
      \sum_{\substack{r,m\ge1\\rm\ge2}}
        \frac{1}{rm}
        \sum_{\pi\in\Irr(G)}
          |c_{r,m,\pi}L_{r,m,\pi}(t)|
      \le
      C_c
      \sum_{k\ge2}\frac{d(k)}{k}\lambda^{1-k}<\infty,
    \]
    with a majorant independent of \(t\).
    \item Therefore, after the finitely many \(rm=1\) terms are reduced by
    adding their explicit logarithmic parts, dominated convergence may be
    passed through the whole fixed-label double determinant series:
    \begin{equation}\label{eq:perron-boundary-fixed-label-majorant-limit}
    \begin{aligned}
      &\lim_{t\uparrow1}
      \left[
        \sum_{r,m\ge1}\frac{\mu(m)}{rm}
          \sum_{\pi\in\Irr(G)}c_{r,m,\pi}L_{r,m,\pi}(t)
        +\log(1-t)\sum_{\pi\in\Irr(G)}c_{1,1,\pi}s_\pi
      \right]                                      \\
      &\qquad =
      \sum_{r,m\ge1}\frac{\mu(m)}{rm}
        \sum_{\pi\in\Irr(G)}c_{r,m,\pi}
          B_{rm,\pi}^{\mathrm{ren}},
    \end{aligned}
    \end{equation}
    where \(B_{1,\pi}^{\mathrm{ren}}\) is the reduced Abel finite part and
    \(B_{k,\pi}^{\mathrm{ren}}=\Log_0\det(I-\lambda^{-k}A_\pi)^{-1}\) for
    \(k\ge2\).
  \end{enumerate}
  No semisimplicity and no strict twisted gap are used here.  The only
  spectral inputs are the algebraic peripheral factorization, the finite
  spectral-radius bound \(\rho(A_\pi)\le\lambda\), and
  \(\lambda>1\).
\end{lemma}

\begin{proof}
  The \(rm=1\) assertion is the determinant factorization in
  \Cref{def:peripheral-peter-weyl-data}.  Namely
  \[
    \det(I-t\lambda^{-1}A_\pi)
    =
    \prod_{\zeta\in\Gamma_\pi}(1-t\zeta)^{m_{\pi,\zeta}}
      Q_{\partial,\pi}(t),
    \qquad Q_{\partial,\pi}(1)\ne0 .
  \]
  The powers \(m_{\pi,\zeta}\) are algebraic multiplicities, so this identity
  is a characteristic-polynomial identity and is unaffected by possible
  nilpotent parts in peripheral Jordan blocks.  Adding \(s_\pi\log(1-t)\)
  cancels exactly the \(\zeta=1\) factors and leaves the Abel branch values
  of the factors \(1-t\zeta\), \(\zeta\ne1\), and of
  \(Q_{\partial,\pi}(t)\).  This gives
  \eqref{eq:perron-boundary-renormalized-b1-factor}.

  For the tail, the finite skew-product block decomposition gives
  \(\rho(A_\pi)\le\lambda\) for every \(\pi\).  Since there are only
  finitely many twisted blocks, there is a constant \(C_0\) such that
  \[
    |\Tr(A_\pi^\ell)|\le C_0\lambda^\ell
    \qquad(\pi\in\Irr(G),\ \ell\ge1).
  \]
  If \(rm\ge2\), then
  \(\rho(t^{rm}\lambda^{-rm}A_\pi)\le\lambda^{1-rm}<1\), uniformly for
  \(0\le t\le1\).  Thus the normalized logarithm is represented by the
  absolutely convergent trace expansion
  \[
    L_{r,m,\pi}(t)
    =
    \sum_{\ell\ge1}
      \frac{t^{rm\ell}\lambda^{-rm\ell}}{\ell}
      \Tr(A_\pi^\ell).
  \]
  The trace bound gives
  \[
    |L_{r,m,\pi}(t)|
    \le
    C_0\sum_{\ell\ge1}\frac{\lambda^{(1-rm)\ell}}{\ell}
    \le C\,\lambda^{1-rm},
  \]
  because \(rm\ge2\) and \(\lambda>1\).  This proves
  \eqref{eq:perron-boundary-fixed-label-majorant-tail}.  Multiplying by the
  bounded coefficients and summing over the finite set \(\Irr(G)\), the
  omitted tail is dominated by
  \[
    C_c
    \sum_{\substack{r,m\ge1\\rm\ge2}}\frac{\lambda^{1-rm}}{rm}
    =
    C_c
    \sum_{k\ge2}\frac{d(k)}{k}\lambda^{1-k},
  \]
  which converges exponentially.  The \(rm=1\) terms are finite in number and
  have the reduced limits already proved.  Dominated convergence for the
  \(rm\ge2\) tail, together with these finite reduced \(rm=1\) limits, proves
  \eqref{eq:perron-boundary-fixed-label-majorant-limit}.
\end{proof}

